% ============================================================
% Round 08: Frontier Map 2026
% ============================================================

\subsection{Round 08：2026 最新科研工作地图}

\subsubsection{基础部分结束}

\begin{summarybox}
\textbf{到这里，基础部分已经打通：}\\
我们已经学完了：
\[
\text{图像编辑任务}
\Rightarrow
\text{Diffusion / Flow 基础}
\Rightarrow
\text{SFT}
\Rightarrow
\text{Reward}
\Rightarrow
\text{Diffusion-DPO}
\Rightarrow
\text{DDPO / DPOK}
\Rightarrow
\text{Flow-GRPO}
\]
后面进入前沿论文阶段。重点不再是“这个公式是什么”，而是读懂每篇新工作到底解决 Flow-GRPO / Diffusion RL 的哪个具体痛点。
\end{summarybox}

\subsubsection{检索说明}

\begin{warnbox}
\textbf{关于 LeapGRPO：}\\
我检索了 \code{LeapGRPO}、\code{Leap-GRPO}、\code{Leap GRPO}、\code{LeapGRPO arXiv}、\code{LeapGRPO GitHub} 等关键词，目前没有找到可靠的 arXiv / GitHub 记录。\\
目前确定存在、且名字非常接近的是 \textbf{LeapAlign}：
\[
\text{LeapAlign: Post-Training Flow Matching Models at Any Generation Step by Building Two-Step Trajectories}
\]
它是 2026-04 的 CVPR 2026 工作，不是 GRPO ratio 路线，而是 two-step leap trajectory + direct reward gradient 路线。后面我们可以单独学它。
\end{warnbox}

\subsubsection{前沿论文总表}

\begin{center}
{\small
\begin{tabular}{p{2.7cm}p{2.1cm}p{4.2cm}p{4.2cm}}
\hline
\textbf{工作（arXiv）} & \textbf{时间} & \textbf{核心问题} & \textbf{关键词} \\
\hline
\href{https://arxiv.org/abs/2505.05470}{Flow-GRPO}
&
2025-05
&
Flow Matching ODE 没有 logprob，如何做在线 RL
&
ODE-to-SDE, GRPO, Denoising Reduction
\\
\href{https://arxiv.org/abs/2508.04324}{TempFlow-GRPO}
&
2025-08
&
Flow-GRPO 对所有 timestep 一视同仁，信用分配太粗
&
trajectory branching, noise-aware weighting, temporal credit
\\
\href{https://arxiv.org/abs/2510.21583}{Chunk-GRPO}
&
2025-10
&
step-level 优化粒度太细，忽略生成过程的时间结构
&
chunk-level policy optimization
\\
\href{https://arxiv.org/abs/2512.21514}{DiverseGRPO}
&
2025-12
&
GRPO 后期容易模式坍缩，生成多样性下降
&
diversity-aware reward, early-stage regularization
\\
\href{https://arxiv.org/abs/2601.00423}{E-GRPO}
&
2026-01
&
低熵步骤探索价值低，但仍消耗大量 RL 更新
&
high-entropy steps, step merging, entropy-aware GRPO
\\
\href{https://arxiv.org/abs/2601.20218}{DenseGRPO}
&
2026-01
&
terminal reward 太稀疏，无法判断每一步贡献
&
dense reward, step-wise reward gain, adaptive stochasticity
\\
\href{https://arxiv.org/abs/2601.05729}{TAGRPO}
&
2026-01
&
I2V 直接套 GRPO 效果不稳定，视频轨迹需要对齐
&
direct trajectory alignment, memory bank, image-to-video
\\
\href{https://arxiv.org/abs/2603.23500}{UniGRPO}
&
2026-03
&
统一模型中同时优化文本推理和图像生成
&
reasoning-driven visual generation, text GRPO + FlowGRPO
\\
\href{https://arxiv.org/abs/2604.04142}{OP-GRPO}
&
2026-04
&
Flow-GRPO 是 on-policy，样本效率低
&
off-policy GRPO, replay buffer, sequence-level IS
\\
\href{https://arxiv.org/abs/2604.06966}{MAR-GRPO}
&
2026-04
&
AR-diffusion 混合模型里 diffusion head 带来 noisy gradients
&
multi-trajectory expectation, uncertainty selection
\\
\href{https://arxiv.org/abs/2604.15311}{LeapAlign}
&
2026-04
&
direct-gradient flow 后训练穿过长轨迹会爆显存和梯度
&
two-step leap trajectory, gradient discounting
\\
\href{https://arxiv.org/abs/2604.18518}{UDM-GRPO}
&
2026-04
&
Uniform Discrete Diffusion 如何稳定接入 GRPO
&
final clean sample as action, forward trajectory reconstruction
\\
\hline
\end{tabular}
}
\end{center}

\subsubsection{二次检索补充：不要漏掉的工作}

为了避免只围绕 Flow-GRPO 这条线看，我又按以下关键词做了二次检索：

\[
\text{GRPO image generation}
\quad
\text{visual generation GRPO}
\quad
\text{image editing GRPO}
\quad
\text{video GRPO}
\quad
\text{autoregressive image GRPO}
\quad
\text{masked image generation RL}
\]

下面这些工作也应该放进我们的前沿阅读地图里。

{\small
\begin{longtable}{p{2.7cm}p{2.1cm}p{4.2cm}p{4.2cm}}
\hline
\textbf{工作（链接）} & \textbf{时间} & \textbf{核心问题} & \textbf{关键词} \\
\hline
\endfirsthead
\hline
\textbf{工作（链接）} & \textbf{时间} & \textbf{核心问题} & \textbf{关键词} \\
\hline
\endhead
\href{https://arxiv.org/abs/2505.07818}{DanceGRPO}
&
2025-05
&
把 GRPO 统一用到 diffusion、rectified flow、T2I、T2V、I2V
&
unified visual GRPO, diffusion + flow, video
\\
\href{https://arxiv.org/abs/2505.17574}{InfLVG}
&
2025-05
&
长视频生成时上下文越来越长，如何用 GRPO 学 context selection
&
long video, inference-time policy, context selection
\\
\href{https://arxiv.org/abs/2505.24875}{ReasonGen-R1}
&
2025-05
&
AR 图像生成先用 CoT 做文本规划，再用 GRPO 优化视觉结果
&
reasoning before image generation, SFT + GRPO
\\
\href{https://arxiv.org/abs/2508.06924}{AR-GRPO}
&
2025-08
&
把 GRPO 直接迁移到 autoregressive image generation
&
AR image generation, token policy, visual rewards
\\
\href{https://arxiv.org/abs/2508.18032}{Visual-CoG}
&
2025-08
&
最终 reward 太粗，给语义推理、过程优化、结果评价三个阶段分别设计 guidance
&
stage-aware RL, chain of guidance
\\
\href{https://arxiv.org/abs/2509.25027}{STAGE}
&
2025-09
&
AR 图像生成里普通 GRPO 会破坏预训练分布，且 token 梯度互相冲突
&
stable AR-GRPO, advantage / KL reweighting, entropy reward
\\
\href{https://arxiv.org/abs/2509.22485}{GCPO}
&
2025-09
&
AR 图像生成中不是所有 image token 都同等重要
&
critical-token policy optimization, AR visual generation
\\
\href{https://arxiv.org/abs/2509.25774}{PCPO}
&
2025-09
&
policy gradient 中不同 timestep credit 不成比例，导致高方差和训练不稳
&
proportionate credit, timestep reweighting
\\
\href{https://arxiv.org/abs/2511.18719}{ViPO}
&
2025-11
&
单个 scalar reward 太粗，忽略图像/视频的空间和时间结构
&
pixel-level advantage, perceptual structuring
\\
\href{https://arxiv.org/abs/2511.19356}{TaRoS}
&
2025-11
&
Video GRPO 中 reward score 会饱和并被当成目标本身，出现 Goodhart 问题
&
target-robust reward signaling, video GRPO
\\
\href{https://arxiv.org/abs/2512.18766}{MaskFocus}
&
2025-12
&
masked generative model 里全轨迹优化成本高，随机挑 step 又不稳定
&
critical steps, information gain, masked generation
\\
\href{https://arxiv.org/abs/2512.13276}{CogniEdit}
&
2025-12
&
细粒度图像编辑只在单步优化，反馈太稀疏，难以控制整条 denoising 轨迹
&
dense GRPO, fine-grained editing, trajectory-level supervision
\\
\href{https://arxiv.org/abs/2512.16864}{RePlan}
&
2025-12
&
复杂指令图像编辑需要先做区域规划，再把规划可靠交给 editor
&
reasoning-guided region planning, planner GRPO, IV-Edit
\\
\href{https://arxiv.org/abs/2601.02256}{VAR RL Done Right}
&
2026-01
&
VAR 生成不同 step 的输入结构异步，普通 GRPO 有 policy conflict
&
asynchronous policy conflict, VAR, mask propagation
\\
\href{https://arxiv.org/abs/2601.03467}{ThinkRL-Edit}
&
2026-01
&
复杂图像编辑需要先推理再生成，不能只靠 denoising 随机性探索
&
reasoning-centric image editing, CoT planning, checklist reward
\\
\href{https://arxiv.org/abs/2602.14186}{UniRef-Image-Edit}
&
2026-02
&
多参考图编辑 / 合成时多个 reference 之间约束冲突
&
multi-reference editing, SELF, MSGRPO
\\
\href{https://www.mdpi.com/1424-8220/26/4/1120}{Flow-Multi}
&
2026-02
&
单 reward 容易 reward hacking，多目标之间会互相冲突
&
multi-reward, Pareto dominance, advantage masking
\\
\href{https://arxiv.org/abs/2603.06623}{GRPO Survey}
&
2026-02
&
Flow-GRPO 之后分支增长太快，需要系统整理方法改进和跨模态扩展
&
survey, credit assignment, reward design, modality extension
\\
\href{https://arxiv.org/abs/2603.09538}{Unified Interleaved GRPO}
&
2026-03
&
统一视觉语言模型要在同一条轨迹里交替生成文本和图像
&
multimodal interleaved generation, hybrid reward, process reward
\\
\href{https://arxiv.org/abs/2603.12648}{MV-GRPO}
&
2026-03
&
单 prompt 单视角评价太稀疏，组内样本关系没有被充分利用
&
augmented condition space, multi-view advantage, dense reward mapping
\\
\href{https://arxiv.org/abs/2603.17461}{AR-CoPO}
&
2026-03
&
streaming AR 视频生成轨迹短、随机性低，SDE-GRPO 不自然
&
contrastive policy optimization, chunk-level neighborhood
\\
\href{https://arxiv.org/abs/2603.26599}{VGGRPO}
&
2026-03
&
视频生成几何一致性差，RGB-space reward 反复 VAE decode 成本高
&
4D latent reward, world consistency, video GRPO
\\
\href{https://arxiv.org/abs/2604.09386}{RC-GRPO-Editing}
&
2026-04
&
图像编辑中全局探索会扰动背景，导致组内 reward 方差变大
&
region-constrained GRPO, localized exploration, attention reward
\\
\href{https://arxiv.org/abs/2604.19234}{OTCA}
&
2026-04
&
多 reward 被压成一个静态标量，并均匀传给整条轨迹
&
objective-aware credit, timestep-aware reward allocation
\\
\href{https://arxiv.org/abs/2604.19406}{HP-Edit}
&
2026-04
&
真实图像编辑缺少可扩展人类偏好数据、scorer 和 benchmark
&
RealPref-50K, HP-Scorer, RealPref-Bench
\\
\hline
\end{longtable}
}

\begin{summarybox}
\textbf{二次检索后的修正：}\\
如果只看 Flow-GRPO、UDM-GRPO、LeapAlign，会漏掉四条很重要的支线：
\[
\begin{aligned}
&
\text{图像编辑专属 RL}
\quad
\text{AR / VAR / masked generation 的 GRPO}
\\
&
\text{视频生成的结构化 reward / credit assignment}
\quad
\text{统一多模态生成的 GRPO}
\end{aligned}
\]
所以后续前沿阅读不应该只沿 Flow Matching 走，而要分成“图像编辑、连续 flow、离散/AR/VAR、视频、统一多模态生成”五条线。
\end{summarybox}

\subsubsection{研究主线总览：从论文列表到研究地图}

这些论文不能按时间顺序硬读。更好的方式是按“它到底在修 Flow-GRPO / Diffusion RL 的哪一个痛点”来分组。

\begin{conceptbox}
\textbf{分类原则：}\\
下面的主线不是互斥分类。一篇论文可能同时属于多个方向，比如 TaRoS 既是 video GRPO，也是在处理 reward hacking；ReasonGen-R1 既是 AR 图像生成，也属于 reasoning-driven generation。这里按“主要研究问题”归档。
\end{conceptbox}

{\small
\begin{longtable}{p{2.45cm}p{3.7cm}p{5.8cm}p{3.3cm}}
\hline
\textbf{研究主线} & \textbf{核心问题} & \textbf{代表工作} & \textbf{读论文时抓什么} \\
\hline
\endfirsthead
\hline
\textbf{研究主线} & \textbf{核心问题} & \textbf{代表工作} & \textbf{读论文时抓什么} \\
\hline
\endhead
连续 flow 的信用分配
&
terminal reward 太粗，不能说明哪个 timestep / chunk / branch 真正导致好结果
&
\href{https://arxiv.org/abs/2505.05470}{Flow-GRPO},
\href{https://arxiv.org/abs/2508.04324}{TempFlow-GRPO},
\href{https://arxiv.org/abs/2510.21583}{Chunk-GRPO},
\href{https://arxiv.org/abs/2601.20218}{DenseGRPO},
\href{https://arxiv.org/abs/2509.25774}{PCPO},
\href{https://arxiv.org/abs/2604.19234}{OTCA},
\href{https://arxiv.org/abs/2603.12648}{MV-GRPO}
&
advantage 是分给 step、chunk、branch、objective，还是多视角 condition
\\
\hline
采样效率、轨迹选择与直接梯度
&
图像 / 视频 rollout 很贵；哪些旧轨迹能复用，哪些步骤值得采样，能不能不走 logprob
&
\href{https://arxiv.org/abs/2604.04142}{OP-GRPO},
\href{https://arxiv.org/abs/2601.00423}{E-GRPO},
\href{https://arxiv.org/abs/2512.18766}{MaskFocus},
\href{https://arxiv.org/abs/2604.15311}{LeapAlign}
&
on-policy / off-policy、replay buffer、entropy step、critical step、direct-gradient 的差别
\\
\hline
reward 设计、多目标与 reward hacking
&
单个 scalar reward 太粗，reward 可能饱和、被钻空子，或者牺牲多样性
&
\href{https://arxiv.org/abs/2512.21514}{DiverseGRPO},
\href{https://www.mdpi.com/1424-8220/26/4/1120}{Flow-Multi},
\href{https://arxiv.org/abs/2511.19356}{TaRoS},
\href{https://arxiv.org/abs/2508.18032}{Visual-CoG},
\href{https://arxiv.org/abs/2511.18719}{ViPO},
\href{https://arxiv.org/abs/2604.19406}{HP-Edit}
&
reward 是单目标、多目标、component-level、pixel-level，还是人类偏好 scorer
\\
\hline
图像编辑专属 RL
&
编辑任务不是单纯生成：既要改对目标区域，又要保护非目标区域，还要能处理复杂指令 / 多参考图
&
\href{https://arxiv.org/abs/2601.03467}{ThinkRL-Edit},
\href{https://arxiv.org/abs/2602.14186}{UniRef-Image-Edit},
\href{https://arxiv.org/abs/2512.13276}{CogniEdit},
\href{https://arxiv.org/abs/2512.16864}{RePlan},
\href{https://arxiv.org/abs/2604.09386}{RC-GRPO-Editing},
\href{https://arxiv.org/abs/2604.19406}{HP-Edit}
&
reward 是否区分 edit region / preserve region；是否先规划再编辑；是否有编辑偏好数据
\\
\hline
离散、AR、VAR、masked / hybrid 生成
&
不是所有视觉生成都是 continuous diffusion / flow；离散 token、VAR、MAR、UDM 的 action 和 logprob 要重新定义
&
\href{https://arxiv.org/abs/2508.06924}{AR-GRPO},
\href{https://arxiv.org/abs/2509.25027}{STAGE},
\href{https://arxiv.org/abs/2509.22485}{GCPO},
\href{https://arxiv.org/abs/2601.02256}{VAR RL Done Right},
\href{https://arxiv.org/abs/2604.06966}{MAR-GRPO},
\href{https://arxiv.org/abs/2604.18518}{UDM-GRPO},
\href{https://arxiv.org/abs/2505.24875}{ReasonGen-R1}
&
action 是 token、clean sample、diffusion head，还是 hybrid trajectory；KL / entropy 怎么稳定
\\
\hline
视频生成与时序一致性
&
视频多了时间维度：trajectory 更长，reward 更容易饱和，运动一致性和跨场景一致性更难
&
\href{https://arxiv.org/abs/2601.05729}{TAGRPO},
\href{https://arxiv.org/abs/2505.07818}{DanceGRPO},
\href{https://arxiv.org/abs/2505.17574}{InfLVG},
\href{https://arxiv.org/abs/2603.17461}{AR-CoPO},
\href{https://arxiv.org/abs/2603.26599}{VGGRPO},
\href{https://arxiv.org/abs/2511.19356}{TaRoS}
&
reward 是否看 motion、geometry、context、trajectory alignment；credit 是否按帧 / chunk 分配
\\
\hline
统一多模态与推理驱动生成
&
模型不再只是生成一张图，而是在同一条轨迹里推理、写文本、生成图像、甚至多轮交错
&
\href{https://arxiv.org/abs/2603.23500}{UniGRPO},
\href{https://arxiv.org/abs/2603.09538}{Unified Interleaved GRPO},
\href{https://arxiv.org/abs/2505.24875}{ReasonGen-R1},
\href{https://arxiv.org/abs/2508.18032}{Visual-CoG}
&
text policy 和 image policy 怎么共存；reward 是 outcome-level 还是 process-level
\\
\hline
综述与全局路线
&
当分支太多时，需要反过来整理整个领域的 taxonomy
&
\href{https://arxiv.org/abs/2603.06623}{Advances in GRPO for Generation Models: A Survey}
&
用它校准分类，不把某一条线误认为全部前沿
\\
\hline
\end{longtable}
}

\begin{summarybox}
\textbf{最适合我们当前学习目标的主线：}\\
因为我们的主题是“图像生成 / 图像编辑侧后训练”，优先级应该是：
\[
\text{图像编辑专属 RL}
\Rightarrow
\text{连续 flow 信用分配}
\Rightarrow
\text{reward 设计}
\Rightarrow
\text{离散 / AR / VAR 扩展}
\Rightarrow
\text{视频和统一多模态}
\]
也就是说，不是先把所有视频、AR、统一多模态论文都读完，而是先围绕图像编辑把核心技术链打穿。
\end{summarybox}

\subsubsection{按目标选择阅读路径}

\begin{enumerate}[leftmargin=1.5em]
  \item \textbf{如果目标是图像编辑：} ThinkRL-Edit $\rightarrow$ RePlan $\rightarrow$ CogniEdit $\rightarrow$ RC-GRPO-Editing $\rightarrow$ HP-Edit $\rightarrow$ UniRef-Image-Edit。
  \item \textbf{如果目标是 Flow-GRPO 算法本身：} Flow-GRPO $\rightarrow$ TempFlow-GRPO $\rightarrow$ Chunk-GRPO $\rightarrow$ DenseGRPO $\rightarrow$ PCPO / OTCA $\rightarrow$ OP-GRPO / E-GRPO。
  \item \textbf{如果目标是 reward 和偏好建模：} DiverseGRPO $\rightarrow$ Flow-Multi $\rightarrow$ Visual-CoG $\rightarrow$ ViPO $\rightarrow$ TaRoS $\rightarrow$ HP-Edit。
  \item \textbf{如果目标是新生成范式：} AR-GRPO $\rightarrow$ STAGE / GCPO $\rightarrow$ VAR RL Done Right $\rightarrow$ UDM-GRPO $\rightarrow$ MAR-GRPO $\rightarrow$ MaskFocus。
  \item \textbf{如果目标是视频：} TAGRPO $\rightarrow$ DanceGRPO $\rightarrow$ InfLVG $\rightarrow$ VGGRPO $\rightarrow$ AR-CoPO $\rightarrow$ TaRoS。
  \item \textbf{如果目标是统一多模态：} ReasonGen-R1 $\rightarrow$ Visual-CoG $\rightarrow$ UniGRPO $\rightarrow$ Unified Interleaved GRPO。
\end{enumerate}

\begin{warnbox}
\textbf{不要被论文名字带偏：}\\
很多论文都叫 GRPO，但它们真正研究的对象不一样。有的研究 \textbf{算法本身}，有的研究 \textbf{reward}，有的研究 \textbf{图像编辑任务约束}，有的只是把 GRPO 迁移到 \textbf{新生成范式}。读之前先判断它属于哪条主线，理解会快很多。
\end{warnbox}

\subsubsection{主线一：Flow-GRPO 的信用分配改进}

Flow-GRPO 的基本做法是：

\[
A_j
=
\frac{
R_j-\text{mean}(R_1,\ldots,R_G)
}{
\text{std}(R_1,\ldots,R_G)+\epsilon
}
\]

然后把同一个 trajectory-level advantage 用到很多 flow steps：

\[
\sum_{k=1}^{K}
\min
\left(
r_{j,k}A_j,
\text{clip}(r_{j,k})A_j
\right)
\]

这个做法能跑通，但问题很明显：

\[
\text{最终 reward 高}
\not\Rightarrow
\text{每一个 timestep 都同样该被奖励}
\]

所以 2025-2026 的一批工作都在解决：

\[
\text{哪个 timestep / chunk / branch 真正导致了 reward 变化}
\]

\textbf{TempFlow-GRPO} 的方向是 temporal credit assignment。它认为 Flow-GRPO 的问题是 temporal uniformity，也就是所有时间步被同样对待。它通过 trajectory branching 把随机性集中到指定 branching point，然后观察最终 reward 差异，从而更精确地归因。

\[
\text{branch at } t_k
\Rightarrow
\text{reward difference}
\Rightarrow
\text{credit for } t_k
\]

\textbf{DenseGRPO} 的方向是 dense reward。它不满足于只用最终 reward，而是用 ODE-based intermediate clean image 去估计每一步带来的 reward gain：

\[
\Delta R_k
\approx
R(\hat{x}_{0\mid t_{k-1}})
-
R(\hat{x}_{0\mid t_k})
\]

这样每一步不再共享同一个 $A_j$，而是有更细粒度的 step-wise reward。

\textbf{Chunk-GRPO} 的方向是 chunk-level credit。它认为单步太细、整条轨迹太粗，所以把连续时间步组成 chunk：

\[
\mathcal{C}_m
=
\{t_a,t_{a-1},\ldots,t_b\}
\]

然后在 chunk 级别优化：

\[
\sum_m
\log p_\theta(\mathcal{C}_m)
A_m
\]

\begin{summarybox}
\textbf{这一类工作的共同问题：}\\
Flow-GRPO 的 terminal reward 太粗。TempFlow-GRPO、DenseGRPO、Chunk-GRPO 都是在回答：
\[
\text{reward 应该分配给哪个时间尺度？step、chunk、branch，还是整条 trajectory？}
\]
\end{summarybox}

\subsubsection{主线二：Flow-GRPO 的采样效率改进}

Flow-GRPO 是 on-policy：

\[
\tau\sim\pi_{\theta_\text{old}}
\]

采样完、更新几轮后，这批样本通常就不能继续用了。

问题是图像 / 视频生成采样很贵：

\[
\text{cost}
\propto
\text{batch size}
\times
\text{group size}
\times
\text{sampling steps}
\]

\textbf{OP-GRPO} 的方向是 off-policy。它把高质量轨迹放进 replay buffer：

\[
\mathcal{B}
\leftarrow
\mathcal{B}
\cup
\{\tau_\text{high-reward}\}
\]

后续训练中重复使用这些轨迹。

但 off-policy 会带来分布偏移：

\[
\tau
\sim
\pi_{\theta_\text{behavior}}
\neq
\pi_{\theta}
\]

所以 OP-GRPO 要做 sequence-level importance sampling correction：

\[
\rho(\tau)
=
\frac{
p_\theta(\tau)
}{
p_{\theta_\text{behavior}}(\tau)
}
\]

并且它观察到 late denoising steps 的 off-policy ratio 容易病态，所以会截断 late steps。

\begin{intubox}
\textbf{直觉：} Flow-GRPO 是“现采现用”；OP-GRPO 是“好轨迹别浪费，放进经验池复用”，但必须处理新旧策略分布不一致的问题。
\end{intubox}

\subsubsection{主线三：探索空间和熵}

Flow-GRPO 通过 SDE 给 ODE 加随机性：

\[
x_{t-\Delta t}
=
\mu_\theta(x_t,t,c)
+
\sigma_t\sqrt{\Delta t}\epsilon
\]

但不是每个时间步都同样需要随机探索。

\textbf{E-GRPO} 的观点是：high entropy steps 更适合 RL 探索，low entropy steps 产生的 rollouts 区分度低，训练信号弱。

所以它会合并连续 low entropy steps：

\[
\text{low entropy steps}
\Rightarrow
\text{merged into one high entropy SDE step}
\]

其他步骤则用 ODE 采样，减少无效随机性。

\textbf{DenseGRPO} 也发现均匀 stochasticity injection 不合适，于是根据 reward-aware signal 调整每个 timestep 的随机性。

\begin{warnbox}
\textbf{这一类工作的核心：} SDE 不是“噪声越多越好”。随机性应该放在真正影响生成结果、且能产生区分度的时间步上。
\end{warnbox}

\subsubsection{主线四：避免 Reward Hacking 和模式坍缩}

GRPO 优化 reward 后，可能出现：

\[
\text{reward 上升}
\quad
\text{diversity 下降}
\]

或者：

\[
\text{metric 高}
\quad
\text{人看起来变单调}
\]

\textbf{DiverseGRPO} 专门处理这个问题。它从两个层面入手：

\begin{itemize}[leftmargin=1.5em]
  \item reward 层面：给语义上更少见的模式 distributional creativity bonus；
  \item generation 层面：在早期 denoising 阶段加入 structure-aware regularization，保护多样性。
\end{itemize}

可以抽象成：

\[
R_\text{total}
=
R_\text{quality}
+
\lambda R_\text{diversity}
-
\beta R_\text{collapse}
\]

它提醒我们：图像后训练不能只看平均 reward，还要看分布层面的 diversity。

\begin{intubox}
\textbf{直觉：} 如果模型发现“某一种构图总能骗到高分”，GRPO 会不断强化这种构图。DiverseGRPO 的目标是让模型既高分，又不要只会一种答案。
\end{intubox}

\subsubsection{主线五：新生成范式上的 GRPO}

前面我们主要学的是 diffusion / flow。2026 的趋势是：

\[
\text{GRPO 不只用于 diffusion / flow}
\]

它正在扩展到统一多模态、离散扩散、AR-diffusion、视频和 3D。

\textbf{UDM-GRPO} 研究的是 Uniform Discrete Diffusion Model。它指出，直接把 GRPO 套到 UDM 上不稳定。它的两个核心设计是：

\begin{itemize}[leftmargin=1.5em]
  \item 把 final clean sample 当作 action，而不是把中间 noisy sample 当作 action；
  \item 用 diffusion forward process 重建 trajectory，使概率路径更贴近预训练分布。
\end{itemize}

抽象对比：

\[
\text{naive UDM-GRPO: intermediate noisy action}
\]

\[
\text{UDM-GRPO: final clean sample as action + forward trajectory reconstruction}
\]

\textbf{MAR-GRPO} 研究的是 masked autoregressive + diffusion head 的混合模型。它发现 diffusion head 会带来 noisy gradients，于是用 multi-trajectory expectation：

\[
\nabla J
\approx
\frac{1}{M}
\sum_{m=1}^{M}
\nabla J^{(m)}
\]

来降低 diffusion trajectory 的随机梯度噪声。

\textbf{UniGRPO} 研究的是统一模型里的 interleaved generation。它同时优化文本推理策略和图像生成策略：

\[
\text{text reasoning GRPO}
+
\text{visual FlowGRPO}
\]

并且为了多轮生成 / 编辑扩展，去掉 CFG 以保持 rollout 线性、用 velocity MSE penalty 替代 latent KL。

\textbf{TAGRPO} 研究 image-to-video。它不只是给整条视频一个 reward，而是用 direct trajectory alignment 在 intermediate latents 上对齐高 reward 轨迹、远离低 reward 轨迹，并用 memory bank 提高效率和多样性。

\begin{summarybox}
\textbf{这一类工作的共同趋势：}\\
GRPO 正在从“Flow Matching 文生图后训练方法”变成更通用的生成模型后训练范式：
\[
\text{continuous flow}
\quad
\text{discrete diffusion}
\quad
\text{AR-diffusion}
\quad
\text{interleaved multimodal}
\quad
\text{video generation}
\]
\end{summarybox}

\subsubsection{主线六：Direct-Gradient Flow 后训练}

并不是所有前沿工作都沿着 GRPO ratio 走。

\textbf{LeapAlign} 是另一条路线：

\[
\text{reward gradient}
\Rightarrow
\text{direct backprop through flow}
\]

它解决的问题是：完整 flow trajectory 太长，直接反传会显存高、梯度爆炸，而且很难更新早期生成步骤。

LeapAlign 把长轨迹压缩成 two-step leap trajectory：

\[
x_k
\Rightarrow
\hat{x}_{j\mid k}
\Rightarrow
x_j
\Rightarrow
\hat{x}_{0\mid j}
\Rightarrow
x_0
\]

其中：

\[
\hat{x}_{j\mid k}
=
x_k-(k-j)v_\theta(x_k,k,c)
\]

\[
\hat{x}_{0\mid j}
=
x_j-jv_\theta(x_j,j,c)
\]

这样 reward gradient 不需要穿过几十步 ODE，只需要穿过两个 leap。

\begin{warnbox}
\textbf{LeapAlign 和 Flow-GRPO 的范式差别：}\\
Flow-GRPO 需要 logprob，适合黑盒 reward。\\
LeapAlign 不需要 logprob，但更依赖可微 reward / surrogate reward，因为它走 direct-gradient。
\end{warnbox}

\subsubsection{建议阅读顺序：以图像编辑为主线}

如果我们后面继续围绕“图像编辑侧后训练”上课，不建议从 UDM-GRPO 或视频论文开始。更合理的顺序是先读图像编辑专属 RL，再回到通用 Flow-GRPO 技术。

\begin{enumerate}[leftmargin=1.5em]
  \item \textbf{ThinkRL-Edit / RePlan：} 先理解复杂编辑为什么需要 reasoning / planning，而不是直接靠 denoising 随机探索。
  \item \textbf{CogniEdit：} 接着看图像编辑里 dense GRPO / trajectory-level supervision 怎么做。
  \item \textbf{RC-GRPO-Editing：} 再看编辑任务最核心的 edit region 与 preserve region 如何分开做 credit assignment。
  \item \textbf{HP-Edit：} 然后看真实图像编辑里的 human preference dataset、scorer 和 benchmark 如何构造。
  \item \textbf{UniRef-Image-Edit：} 最后扩到多参考图编辑，理解 reference conflict 和 multi-source reward。
  \item \textbf{DenseGRPO / TempFlow-GRPO / Chunk-GRPO：} 回头补通用 Flow-GRPO 信用分配，为理解上面几篇提供算法底座。
  \item \textbf{OP-GRPO / E-GRPO / LeapAlign：} 再看效率、采样步骤选择，以及不用 logprob 的 direct-gradient 路线。
  \item \textbf{AR-GRPO / UDM-GRPO / MAR-GRPO：} 最后再进入离散、AR、hybrid 生成范式。
\end{enumerate}

\begin{summarybox}
\textbf{下一节建议：图像编辑专属 RL}\\
我们现在最该学的不是“又一个 GRPO 变体”，而是：
\[
\text{当任务从 text-to-image 变成 image editing 时，reward、action、trajectory、credit assignment 分别要怎么改？}
\]
所以下一节可以从 ThinkRL-Edit、RePlan、CogniEdit、RC-GRPO-Editing 这条线开始。
\end{summarybox}

\subsubsection{参考链接}

\begin{itemize}[leftmargin=1.5em]
  \item \textbf{Flow-GRPO 基础与信用分配：}
  \href{https://arxiv.org/abs/2505.05470}{Flow-GRPO},
  \href{https://arxiv.org/abs/2508.04324}{TempFlow-GRPO},
  \href{https://arxiv.org/abs/2510.21583}{Chunk-GRPO},
  \href{https://arxiv.org/abs/2601.20218}{DenseGRPO},
  \href{https://arxiv.org/abs/2509.25774}{PCPO},
  \href{https://arxiv.org/abs/2604.19234}{OTCA},
  \href{https://arxiv.org/abs/2603.12648}{MV-GRPO}。
  \item \textbf{效率、探索与直接梯度：}
  \href{https://arxiv.org/abs/2604.04142}{OP-GRPO},
  \href{https://arxiv.org/abs/2601.00423}{E-GRPO},
  \href{https://arxiv.org/abs/2512.18766}{MaskFocus},
  \href{https://arxiv.org/abs/2604.15311}{LeapAlign}。
  \item \textbf{reward 与偏好建模：}
  \href{https://arxiv.org/abs/2512.21514}{DiverseGRPO},
  \href{https://www.mdpi.com/1424-8220/26/4/1120}{Flow-Multi},
  \href{https://arxiv.org/abs/2511.19356}{TaRoS},
  \href{https://arxiv.org/abs/2508.18032}{Visual-CoG},
  \href{https://arxiv.org/abs/2511.18719}{ViPO},
  \href{https://arxiv.org/abs/2604.19406}{HP-Edit}。
  \item \textbf{图像编辑专属：}
  \href{https://arxiv.org/abs/2601.03467}{ThinkRL-Edit},
  \href{https://arxiv.org/abs/2512.16864}{RePlan},
  \href{https://arxiv.org/abs/2512.13276}{CogniEdit},
  \href{https://arxiv.org/abs/2604.09386}{RC-GRPO-Editing},
  \href{https://arxiv.org/abs/2604.19406}{HP-Edit},
  \href{https://arxiv.org/abs/2602.14186}{UniRef-Image-Edit}。
  \item \textbf{离散、AR、VAR、hybrid：}
  \href{https://arxiv.org/abs/2508.06924}{AR-GRPO},
  \href{https://arxiv.org/abs/2509.25027}{STAGE},
  \href{https://arxiv.org/abs/2509.22485}{GCPO},
  \href{https://arxiv.org/abs/2601.02256}{VAR RL Done Right},
  \href{https://arxiv.org/abs/2604.06966}{MAR-GRPO},
  \href{https://arxiv.org/abs/2604.18518}{UDM-GRPO}。
  \item \textbf{视频与统一多模态：}
  \href{https://arxiv.org/abs/2601.05729}{TAGRPO},
  \href{https://arxiv.org/abs/2505.07818}{DanceGRPO},
  \href{https://arxiv.org/abs/2505.17574}{InfLVG},
  \href{https://arxiv.org/abs/2603.17461}{AR-CoPO},
  \href{https://arxiv.org/abs/2603.26599}{VGGRPO},
  \href{https://arxiv.org/abs/2603.23500}{UniGRPO},
  \href{https://arxiv.org/abs/2603.09538}{Unified Interleaved GRPO},
  \href{https://arxiv.org/abs/2505.24875}{ReasonGen-R1}。
  \item \textbf{综述：}
  \href{https://arxiv.org/abs/2603.06623}{Advances in GRPO for Generation Models: A Survey}。
\end{itemize}
